
%=====================================================================
\textit{This chapter evaluates the finished project and reflects on the technical, methodological, and contextual factors that shaped the outcome.}

\section{Overview of End Product}

The finished product meets the main requirement the customer requested, with few features from the original project description being omitted. The quality of the finished project as a concept test for new visualization methods was noted in the customer feedback presented in Section~\ref{sec:customer_feedback}.

The overall solution was assessed as satisfactory by the customer, as it allowed them to test concepts of new visualization modes like signal strength, clustering, and heat map layers. Collaboration with the customer was effective troughout the project, with frequent communication supporting the addition of new features and the resolution of technical questions as they arose.

The project demonstrates that visualization of large datasets through methods not currently in use by the customer can provide concrete value. This is reflected in the customer's stated intention to incorporate several of these features into their existing software.

\section{Discussion of Technical Solution}

The technical solution, described in Section~\ref{sec:system_architecture}, is documented and functions as intended, with emphasis on simplification through the use of helper classes and repeated refactoring throughout the project.


\subsection{System Architecture}

The system architecture, with Deck.gl overlaid on MapLibre as the base map and a Docker-based backend providing data, worked well throughout development. The modular approach led to MapLibre being separated from the Deck.gl implementation (Section~\ref{sec:deckgl}), with the map itself responsible only for  map-related functionality such as zooming and movement. This separation simplified the overall complexity and produced cleaner code with a clear separation of concerns. Achieving this required a period of refactoring tightly coupled code, particularly where the map itself contained functionality unrelated to its task. 

The architecture choice was partially guided by the customer, whose request to use React and other technologies already in use within their infrastructure shaped the group's decisions, alongside with the customer providing the backend. Maplibre was suggested in the project document and was ultimately selected as the map provider, supported by its open-source nature and use of vector tiles (Section~\ref{sec:vector tiles}).

\subsection{Frontend}

The frontend implementation successfully followed the wireframes established at the beginning of the project (Section~\ref{sec:design_process}). The design is minimalist by comparison to other applications, prioritizing the map and visualization components over additional interface elements. 

The choice of map provider had both advantages and disadvantages. Maplibre is a fork of the now closed-source Mapbox, which meant that documentation was less comprehensive than required, and that more detailed information for equivalent code often had to be retrived from the Mapbox website. As a result working with MapLibre required more time spent understanding the library than was originally anticipated, delaying implementation beyond what had been planned.

Deck.gl (Section~\ref{sec:deckgl}) proved to be an effective choice for visualization, providing all functionality needed to implement the customer's requested features. Layering through Deck.gl allowed straightforward integration with MapLibre by utilizing a MapBox-style overlay. Since MapLibre is a fork of MapBox, no additional integration code was required beyond what would have been needed for MapBox itself. 

% Needs to be fully rewritten/ add more context since we have experience, just not the knowledge needed for developing a map based application and the logic that should be used as well as standards.
The use of React with TypeScript enforced stronger type safety than plain React, but introduced an initial learning curve. The group's lack of experience with the coding standards needed for react and typescript, led to development time being spent on learning the basics as well as more advanced tools such as Context, which was ultimately used to simplify the code.  

The single-page design, which covered all required functionality through interactive components, suited the project well. One advantage over a multi-page architecture, was that the map and visualization remained visible at all times, including during manual testing of other components. The overall simplicity of the interface also aligned with the projects focus on concept testing, particularly as the customer themselves had requested that less time would be spent on design and more on functionality.
 

Overall, the frontend implementation fulfilled the project's main goals while highlighting how the use of third party tools can present challenges, and how tool selection is an important step in any software development project. 


\subsection{Visualization Methods and User Interaction}

The chosen visualization methods were effective for different aspects of anomaly visualization, each with their own strenghts and weaknesses related to readability, performance and interaction. 

% TODO: Add the date for change in efficeny for clustering and singular scatter plot.
Clustering, implemented as described in Section~\ref{sec:clustering_method} and shown in Section~\ref{sec:effect_of_clustering_results}, produced a measurable performance improvement compared to a normal scatter plot of all points in the dataset, while also improving readability. 

The zoom function with layer recalculation was observed to remain responsive during manual testing, although responsiveness remained dependent on the hardware running the application. The layer was implemented with a focus on simple visualization, including a numeric label indicating count of anomaly groups within each cluster, presented in a clear and readable manner. The layer is also interactive, with pressing 
% Might need to be rewritten cause it is easily misunderstood.
a single individual anomaly group cluster leading to the anomaly group layer, this is performant through api call to the backend, instead of acquiring the information through frontend filtering. The main drawback with clustering is the lack of specific details of each anomaly group, requiring multiple layer changes to go into an anomaly group and then into the individual anomalies, it gives a generally simpler overview of where anomalies are grouped, and how numerous they are, but no specific details on the anomalies inside. Although this is worth the tradeoff versus having a normal scatterplot with a chaotic amount of points shown at once. 

The Heat map visualization, implemented as described in Section~\ref{sec:Heat_Map_method}, provides value by representing the entire dataset, rather than the aggregated form used by clustering. This produces a overview of all data points based on proximity, in contrast to clustering visualization, it utilizes color intensity to represent hot spots more precisely. Color based encoding is more visually intuitive for identifying hot spots compared to visually identifying density based on size and numeric labels. This benefit also justifies higher rendering costs utilized by heat map's. This means the visualisation method is slower than other methods while it offers a more comprehensive overview of the dataset. % TODO: add data here or in results for rendering times like for clustering and induvidual scatter plot.

The anomaly group layer, described in Section~\ref{sec:anomalies and anomaly groups}, provides detailed information about each individual group. The signal strength visualization is central to this layer, and the implementation received positive feedback from the customer (Section~\ref{sec:customer_feedback}). The combination of location, anomaly type, and connections to base stations and satellites in the information component,provides a clear overview of how an anomaly group relates to the receivers of AIS messages. 

The individual anomaly layer, implemented as a scatter plot (Section~\ref{sec:scatter_plot}), was added at the customer's request, with an info component being used to display the data properties of each anomaly. The layer provides further insight into the individual anomalies within a group, including their heading and signal strength. The layer adds clear value to the project and presents no significant drawbacks. 

The remaining visualizations, including the base station layer (Section~\ref{sec:base_station}) and the satellite layer, are required for layers such as the anomaly group with signal strength to provide any useful information. They are therefore crucial too the project and presents no notable drawbacks in their implementation. 

User interaction is a core part of the application, and user's should be able to navigate and retrieve information without noticeable slowdown or errors. Based on manual testing during development, user interaction was observed to remain responsive, even though performance remains heavily dependant on hardware used, with the available GPU acting as the main limiting factor for rendering performance (Section~\ref{sec:hardware_acceleration}).% TODO: Add measured data for interaction latency

Actual user interaction is based heavily around the zoom and camera location of the application. This interaction with the map has drawbacks such as performance loss, but allows for information on anomalies, filtering and searching. This user interactivity comes with the downside of the performance loss, with system responsiveness being affected. As a core feature of the application, the trade of with responsiveness and performance are offset, with non static datasets being much easier to analyse. 


\subsection{Backend and data handling}

The backend being provided by the customer, has comprehensive documentation of endpoints through swagger, with all the endpoints needed for the project implemented. Docker proved to simplify the process of running the backend, with few sentences of code being required to start and fill the database with test data. Observing the backend performance showed no critical performance or data errors, but the delivered backend itself comse with limitations on being constrained to the predefined API structure, with extensions being more complicated for the group to implement themselves, due to lack of experience with the programming language GO. Although the observed performance was satisfactory, testing of the performance was not controlled by the group itself, with core negatives of the backend implementation being wait time for customer to implement required features, a less often occurrence. As a result the backend was well implemented, but not controlled solely by the group, this being a tradeoff.

Data handling through GeoJSON, proved a good choice for both the backend and frontend, with the libraries used already supporting the data. The data itself being parsed through a deck.gl layer. The layer being responsible for parsing data from each GeoJSON feature. Proved to be an effective design, that allows for simple access to data, through accessing properties and geometry of said features. A tradeoff of using GeoJSON is that the frontend must have prior knowledge of the data structure in order to correctly access said properties and geometry. This introduces dependency on a clear schema. Another issue with data, has been asynchronous errors, with wrong data being delivered, or data being accessed before delivery. 
  
%=====================================================================

\section{Performance and Scalability Discussion}

The performance of the system is highly dependent on hardware specifications, with particular effect coming from GPU capability. Webgl has highest effect on greater GPU's, leading to performance being bottlenecked on weaker systems. 


\subsection{Rendering Performance}

Base rendering performance of maplibre, is observed to not have issues, with performance being taxed on the visualization methods and dataset size, rather than the map. Datasets were tested from thousands, to over one hundred thousand, with both having little issues on group computers. The architectural design of the project is not bottlenecked based on the backend, but with frontend using supercluster to recalculate clusters based on zoom, and similar calculations being done on heat map, leads to zoom and move operations being most taxing on the system. This can be tweaked to recalculate fewer times on zoom level, but will lead to a less responsive system. The core bottleneck of any performance is then the frontend itself, with the backend being called in milliseconds. 


%=====================================================================

\section{Project Execution}

This section evaluates the project management process described in Section \ref{sec:dev_methodology}, focusing on how the chosen methodology, team structure, and communication practices influenced the development outcome.

Although the adopted methods ensured a structured development process, several of them introduced weekly overhead that proved disproportionate for a two-person team operating under a fixed deadline. The following subsections evaluate each aspect individually.

\subsection{Scrum Framework and Sprint Structure}

The Scrum framework provided a consistent structure for planning and delivery. The system was designed for large groups, so it had issues in relation to workload and redundancy for groups of two. Therefore, the daily stand-up system was altered and implemented to occur multiple times a day instead to make sure the group had proper accountability and focus in relation to scope. It also helped keep motivation high and allowed the group to quickly make changes to tasks depending on problems found with them. Instead of the standard two-week sprint, the group used weekly sprints as well, to supplement issues which were fixed by multiple daily stand-ups. Since it was a group of two, the focus of the project and scope were defined and worked toward each week while also making sure the backlog was handled.

However, several Scrum artifacts introduced documentation overhead that provided limited value at this team size. This became a bigger issue because of the decision to have weekly sprints, which increased the overhead that needed to be handled compared to the standard system. The problem was amplified because sprints documented and reviewed the backlog every week, even though group members had continuous visibility into task status through GitHub Projects.
 % Maybe ref or  section

In relation to sprint documentation, the sprint artefacts provided traceable project documentation, which served as a reference when writing the report while ensuring important information was not lost in miscommunication.

\subsection{Team Composition and Workload Distribution}

In a group of two where both members had mostly the same knowledge and expertise, workload was distributed based on member preference and the nature of each task, which made documentation responsibilities formally split where one member handled the sprint documentation while the other handled meeting documentation while development tasks were shared. In practice however, more demanding tasks were rotated between members to maintain motivation. This type of informal balancing worked well for a group of 2 but would not scale to a larger team, where responsibilities would be split more evenly with explicit role definitions. 

\subsection{Communication and Meeting Structure}

Slack provided an efficient channel for direct communication with the customer, enabling rapid resolution of short technical questions. The customer's domain expertise allowed complex questions to be resolved within short discussions rather than through a long discussion. While communication with the supervisor was handled primarily through email, with in person consultations used for urgent issues, because most work was conducted on campus.

For meetings, biweekly meetings were held with both the supervisor and the customer at Kystverket. Before these meetings, discussions were held to discuss which questions should be asked. However, most of these discussions happened through a small meeting before the actual meeting or through daily stand-ups within the group.

Most meetings with the customer were focused on whether the customer liked what had been developed, followed by technical questions to figure out the difficulty of harder problems and what they wanted the group to focus on. With the supervisor, however, problems relating to documentation and any issues with what had been done were mostly discussed. There were sometimes occasions where meeting times were delayed because the customer was busy or had taken a break from work.

\subsection{Task Management and Version Control}

The task management infrastructure described in Section \ref{sec:task_management} worked effectively for the project, although the multi-repository setup introduced some context-switching overhead. The four workflow stages would have been more useful in a larger group, however since work was handled in person they served mainly to track which tasks had not been done or which member was currently working on each task. The three-repository structure occasionally caused issues to be filed in the wrong board, where they were overlooked until rediscovered and relocated. The burndown chart was generated automatically and used to track project progress and identify whether the project was accumulating unaddressed issues. The three-board structure worked better than a single combined board, because GitHub Projects allowed all three to be viewed together at the overview level while keeping each repository's issues separated at the detail level.

For version control, Git was utilized because it allows code to be reverted easily if new code breaks the system or creates unexpected issues. Branching was useful overall, even though it occasionally caused merge conflicts and required refactoring after code was combined. This was still preferable to working in a single shared branch, where both members would have caused conflicts continuously during development.

\subsection{Scope Management}

Scope management was handled reactively through customer meetings rather than through a formal change process, which had both benefits and drawbacks. The weekly sprint cycle allowed scope changes to be incorporated quickly, which was valuable given the short development period and limited team capacity. This approach kept development aligned with the customer's interests and prevented time from being spent on challenging features that were not needed. However, the absence of a formal scope freeze meant that new requests led to planned features being cut to accommodate them. The reactive approach worked well within the constrained timeline, although a longer project would have permitted more features to be properly developed and polished. For a project of similar length, a formal scope freeze date would improve predictability. Weighted prioritization could also help larger teams manage competing requests, but the small team size and direct communication with the customer made this less necessary in this project.

% Should be in technical solution and discussion of that since 
\begin{comment}
\subsection{Code Quality and Maintainability}

The code quality was designed to follow software principles, but due to lack of experience with react itself, some uncertainties about use of specific methods arose. The use of prior programming experience to create files that handle specific tasks, are however used to simplify the program, with refactoring being done multiple times to follow these guidelines

\subsection{problems during development}

Multiple problems arose during the development of the project. A core problem being the changing of the scope throughout the project, which lead to features being cut in focus on new features being added.This was handled through communication with the customer, where functionality to focus on was honed in on. 
\end{comment}

%=====================================================================

\section{Limitations}

Although the project fulfilled the main end requirements through scope changes that the the customer requested, it still has limitations and is a concept testing software. 

\subsection{Technical Limitations}

Technical limitations include not having implemented real AIS data from the customer in the project. This would be unlikely to have happened with any amount of work on the project as such information is classified, resulting in it being something the customer can take forward themselves. Other limitations are satellite orbit not being plotted, as was a request that was given, but not able to implement in time. This would have allowed for users to see the orbit from the satellite, from its current position, to its positions in the future. Although satellite orbit was not implemented, satellites themselves were, and show up in the signal strength visualization. Another feature from the project outline not implemented, that was quietly removed from scope, is automatic report generation, but this was removed as a request in the project, where focus went into concept testing. 




%=====================================================================